\documentclass[10pt,a4paper]{amsart}
\usepackage[margin=19mm]{geometry}
\usepackage{amsmath,amssymb,mathtools}
\usepackage[scr=rsfs,cal=euler]{mathalfa}
\usepackage{xfrac}
\usepackage[unicode,hidelinks,bookmarks=false]{hyperref}
\newcommand{\ie}{i.e.\ }
\newcommand{\cf}{cf.\ }
\newcommand{\resp}{respectively\ }
\newcommand{\Subsection}{\subsection}
\providecommand{\nobreakdash}{\nobreak\hbox{-}}
\setlength{\emergencystretch}{2em}
\multlinegap0pt
\usepackage{tikz}
\usepackage{mathtools}
\usepackage{amssymb}
\usepackage{xcolor}
\newtheorem{theorem}{Theorem}
\newtheorem{lemma}{Lemma}
\newcommand{\Ext}{\operatorname{Ext}}
\newcommand{\Hom}{\operatorname{Hom}}
\newcommand{\Int}{\operatorname{Int}}
\def\TT{\mathbf{T}}
\renewcommand{\k}{\mathbf{k}}
 \renewcommand{\mod}{\operatorname{mod}}
\def\surf{\mathbf{S}}                       %FST's surface
\newcommand{\udim}{\operatorname{\underline{dim}}}
\title{A geometric model for the non-$\tau$-rigid modules of type $\widetilde{D}_n$}
\author{Blake Jackson}
\date{Source: arXiv:2507.02218v2 (CC BY 4.0)}
\begin{document}
\maketitle

\noindent\emph{Source and licence.} A geometric model for the non-$\tau$-rigid modules of type $\widetilde{D}_n$. Version arXiv:2507.02218v2 (CC BY 4.0), distributed by arXiv under the Creative Commons Attribution 4.0 International licence, http://creativecommons.org/licenses/by/4.0/, as recorded in the arXiv licence metadata for this exact version and shown on the abstract page https://arxiv.org/abs/2507.02218v2. Full licence notice: \texttt{licenses/arxiv-2507.02218-CC-BY-4.0.txt}.\par\smallskip

\noindent\textit{Editorial scope.} Counted unit: the article's lemma lemma (source0.tex lines 1340--1343). The article's complete original proof of it is delivered with it (source0.tex lines 1345--1363, one \texttt{proof} environment of 19 lines). No mathematics is written, repaired or re-derived: the statement and the proof are the article's own lines, in the article's own notation, and no retained line is substituted or re-flowed.  Retained uncounted as the source-local context the proof needs: lines 466--469.  Nothing was omitted from any retained span, so the package carries no omission record.  No retained line contains a citation, so the wrapper carries no bibliography and no bibliography entry.  No retained line contains a figure environment, an image-inclusion command or drawing code, so no image is delivered and the package declares no asset dependency.  Editorial formatting only, in addition: an AMS \texttt{amsart} preamble that restates the article's own declarations and macros used by the retained lines, namely the environments \texttt{theorem}, \texttt{lemma} on separate counters, together with the standard packages those lines need.  The retained environments print in this document as Theorem 1, Lemma 1 in order of appearance, whereas the article prints the same items with the numbering of its own full document.  The complete native source of the article as distributed in the arXiv e-print for this exact version is bundled unchanged as \texttt{sources/180931/source0.tex}.
\begin{theorem}\label{thm-0-homs}
    Let $\Gamma(\mod \k Q )$ be the Auslander-Reiten quiver of the module category of the path algebra over an acyclic quiver of type $\widetilde{D}_n$. 
    Then \[ \Hom (\mathcal{Q}(\k Q) ,\mathcal{P}(\k Q)) = \Hom (\mathcal{R}(\k Q) ,\mathcal{P}(\k Q)) = \Hom (\mathcal{Q}(\k Q) ,\mathcal{R}(\k Q)) = 0. \]
\end{theorem}

\begin{lemma}
    Let $(\gamma_1, \kappa_1, -\infty) \in \mathcal{P}(\TT)$ and $(\gamma_2, \kappa_2, \lambda) \in E^\times_\lambda$ where $\TT$ is an acyclic triangulation of a twice-punctured $(n-2)$-gon $\surf$ and let $M(\gamma_1, \kappa_1, -\infty) = M_1 \in \mathcal{P}(\k Q^\TT)$ and $M(\gamma_2, \kappa_2, \lambda_2) = M_2 \in \Gamma(\mod \k Q^\TT)$ be the corresponding modules under the equivalence $\varphi^{-1}$.
    Then $$\Int( (\gamma_1,\kappa_1, -\infty), (\gamma_2,\kappa_2, \lambda) ) = \dim_\k \Ext^1 (M_1, M_2) + \dim_\k \Ext^1 (M_2, M_1).$$
\end{lemma}

\begin{proof}
    By Theorem~\ref{thm-0-homs}, we know that any nonzero homomorphism between $M_1$ and $M_2$ must begin in the preprojective component.
    If $M_2 \in \mathcal{P}(\k Q^\TT)$ as well, assume without loss of generality that $M_1$ is to the left of $M_2$ in the Auslander-Reiten quiver.
    First note that, for an indecomposable modules $M$, we have \[\Hom (P(j), M) \cong M_j\] where $M_j$ is the vector space at vertex $j$ in the indecomposable quiver representation $M$.
    Since $\Int( -, - )$, $\Hom(-,-)$, and $\Ext^1 (-, -)$ are invariant under $\tau$ and $\rho$ for non-projective modules over a path algebra of type $\widetilde{D}_n$, we can apply these translations to both objects until $(\gamma_1, \kappa_1, -\infty) \in \TT$.
    Say $\rho^{n+1} (\gamma_1,\kappa_1, -\infty) = j \in \TT$ so that $\tau^{n} M_1 = P(j)$.
    Then, if $e_j$ is the unit vector with a $1$ in the $j^{th}$ position and zeros elsewhere and $\tau^{n+1} M_2 \neq 0$,
    \begin{align*}
        \Int( (\gamma_1,\kappa_1, -\infty), (\gamma_2,\kappa_2, \lambda) ) &= \Int( \rho^{n+1}(\gamma_1,\kappa_1, -\infty), \rho^{n+1}(\gamma_2,\kappa_2, \lambda) ) \\
        &= \Int( j, \rho^{n+1}(\gamma_2,\kappa_2, \lambda) ) \\
        &= \udim (\tau^{n+1} M_2) \cdot e_j = (\tau^{n+1} M_2)_j \\
        &= \dim_\k \Hom (P(j), \tau^{n+1} M_2) \\
        &= \dim_\k \Hom (\tau^{-n}P(j), \tau^{-n}\tau^{n+1} M_2) \\
        &= \dim_\k \Hom (M_1, \tau M_2) \\
        &= \dim_\k \Ext^1 (M_2, M_1)
    \end{align*}
    where the last inequality follows from the Auslander-Reiten formulas.
    Moreover, by our assumption that $M_1$ is to the left of $M_2$, we have that $\Ext^1 (M_1, M_2) = 0$.
\end{proof}

\end{document}
